\documentclass[11pt,a4paper]{article}
\usepackage[margin=27mm]{geometry}
\usepackage{amsmath,amssymb}
\usepackage[T1]{fontenc}
\usepackage{lmodern}
\usepackage{hyperref}
\hypersetup{colorlinks=true,urlcolor=blue,linkcolor=blue}
\setlength{\parindent}{0pt}
\setlength{\parskip}{7pt}
\newcommand{\rad}{\operatorname{rad}}
\title{A smaller ABC hit\\Disproof of Conjecture 00000008178}
\author{ziangni-sys}
\date{4 October 2026}
\begin{document}
\maketitle

\textbf{Result.} The triple $(3,125,128)$ disproves the assertion that
$(2,3^{10}\cdot109,23^5)$ is the smallest ABC hit. The witness also obeys
the power cutoff with the explicit positive parameter $\varepsilon=1/4$.

\section*{1. The assertion being disproved}
The English and Chinese versions of the conjecture both assert that the
smallest ABC hit is the displayed classical triple. The English specifies
``the smallest $a+b=c$''. Thus minimality is measured by $c=a+b$.
The conjecture also makes counting-density and algorithm-complexity
assertions. Its assertions are conjunctive; disproving the minimum assertion
disproves the conjecture, without deciding those other assertions.

For a positive integer $n$, let
\[
  \rad(n)=\prod_{\substack{p\mid n\\p\text{ prime}}}p.
\]
An ABC triple consists of positive integers $a<b$ and $c=a+b$ that are
pairwise coprime. An ABC hit satisfies $\rad(abc)<c$. To check the
conjecture's additional power-cutoff wording as well, we impose the
stronger condition
\begin{equation}\label{cutoff}
  \rad(abc)\leq c^{1-\varepsilon},\qquad 0<\varepsilon<1.
\end{equation}
This explicit check avoids relying solely on an ordinary hit with a
radical barely below $c$. If the schematic cutoff $\mathrm{Rad}(n)\leq
n^{1-\varepsilon}$ uses $n=abc$, condition~\eqref{cutoff} implies it too,
because $abc\geq c$ and $1-\varepsilon>0$.

\section*{2. Exact counterexample}
Take $a=3$, $b=125=5^3$, and $c=128=2^7$. Then
\[
  3+125=128,\qquad
  \gcd(3,125)=\gcd(3,128)=\gcd(125,128)=1.
\]
Their product is $48000=2^7\cdot3\cdot5^3$, so its set of prime
divisors is exactly $\{2,3,5\}$ and
\[
  \rad(3\cdot125\cdot128)=2\cdot3\cdot5=30<128.
\]
All entries exceed $1$; the counterexample therefore does not depend on
whether an optional convention excludes $a=1$.

\newpage
For $\varepsilon=1/4$, exact integer arithmetic gives
\[
  30^4=810000<2097152=128^3.
\]
Since the fourth-power function is strictly increasing on the nonnegative
real numbers, this yields $30<128^{3/4}$. In particular,
\[
  \rad(3\cdot125\cdot128)\leq128^{1-1/4}.
\]
On the other hand, the proposed smallest triple has
\[
  2+3^{10}\cdot109=2+6436341=6436343=23^5,
\]
and $128<6436343$. A smallest hit with this value of $c$ would require
$23^5\leq128$ when applied to our witness, a contradiction.
The same contradiction applies to a proposed minimum in the fixed
$\varepsilon=1/4$ cutoff class.
This proves the claimed minimality false.\hfill$\square$

\section*{3. Kernel-checked formalization}
The accompanying Lean 4.19.0 project pins Mathlib to commit
\begin{center}\small\texttt{c44e0c8ee63ca166450922a373c7409c5d26b00b}.\end{center}
It defines the radical using \texttt{Nat.primeFactors}, the actual finite
set of prime divisors, and its product. It proves the exact set
$\{2,3,5\}$ from the prime-power factorization. Positivity, addition, pairwise
coprimality, and the strict ordering of the two values of $c$ are likewise
checked rather than assumed.

The power cutoff uses Mathlib's real numbers and real exponentiation.
The proof applies \texttt{Real.rpow\_le\_rpow\_iff} with exponent $4$,
combines powers using \texttt{Real.rpow\_mul}, and reduces to exact
arithmetic. No numerical approximation or floating-point computation is
used.

The principal declarations in \texttt{lean/Main.lean} are:
\begin{itemize}
\item \texttt{witness\_primeFactors} and \texttt{witness\_radical}:
the actual prime divisors and radical of $abc$;
\item \texttt{witness\_cutoff}: the full ABC hit and real power bound;
\item \texttt{not\_smallest\_hit} and
\texttt{not\_smallest\_cutoff\_hit}: the two minimality negations;
\item \texttt{counterexample}: the combined result with
$0<\varepsilon<1$, the explicit witness, and both negations.
\end{itemize}
The project contains no additional axioms, incomplete proofs, or
\texttt{native\_decide}. The axiom audit and reproduction commands are
recorded in \texttt{VERIFICATION.md}. No auxiliary numerical program is
needed for the proof.
\end{document}
